\chapter{Introduction}
\label{sec:introduction}
\label{sec:motivation}

Answer Set Programming (ASP) is a declarative approach to combinatorial search. Instead of describing an algorithm step by step, the modeller describes the properties that an admissible solution must satisfy, and a solver enumerates the assignments that satisfy them \cite{gelfond1988stable, schaub2024aspPractice, clingoGuide}.

ASP solvers are essentially highly optimized backtracking, or more appropriately \emph{backjumping}, engines. Modern ASP solving is dominated by two main paradigms: ground-and-solve and lazy grounding.

In ground-and-solve systems, such as clingo, first-order rules are instantiated over the problem domain before search starts, including the domain-specific heuristics that direct the solver's search. 
This approach can be severely memory-intensive, with grounding bottlenecks manifesting in two main scenarios: as problem instances scale, and whenever domain-specific heuristics depend on dynamic search state, for instance through an aggregate over the atoms that are true \emph{at that point of the search}. 
Because grounding precedes search, the grounder must enumerate every possible value the aggregate could take, and a single heuristic ends up producing a number of ground objects that is combinatorial in the size of the instance. In contrast, lazy-grounding systems interleave grounding and solving. This can reduce memory use and extend reach on memory-intensive problems, but interleaving grounding and solving can also make search slower.

Depending on its underlying architecture, an ASP solver can be seen as an engine that compiles or interprets a very convenient and elegant declarative language \cite{calimeri2020aspcore2} to search for a solution. The convenience of that language is paid for in control: the modeller states \emph{what} an admissible solution is, but not \emph{how} the engine should look for one. A pure backtracking approach behaves as an uninformed search, and on problems whose structure is known by construction, guiding the solver with that knowledge is extremely helpful.

\emph{Domain-specific heuristics} are the means by which the modeller feeds that knowledge into the search. Syntactically they are declarative annotations that the grounder instantiates like ordinary rules, and their purpose is to tell the solver which atom to decide next, and with which polarity; their impact on solving is well documented \cite{gebser2013heuristics,dodaro2016combining}. They are, however, part of the program, and clingo instantiates the program before search starts. A heuristic is therefore expanded, in full, before the first decision is taken.
That expansion can be relatively harmless as long as the heuristic depends only on the input. It stops being harmless as soon as the heuristic depends on \emph{how the search is going}, which is precisely when a heuristic is most useful. Consider the directive used throughout this thesis to guide the Balanced Sum Problem (BSP), the problem of splitting $\{1,\ldots,n\}$ into two sets \texttt{b/1} and \texttt{c/1} of nearly equal sum:

\begin{minted}{prolog}
#heuristic b(X) : x(X), not c(X), S = #sum{Y : c(Y)}, W = ((n+1)*S)+X. [W, true]
\end{minted}

The rule gives a higher weight to placing $X$ in $b$ as the current sum $S$ in $c$ increases.
The aggregate value $S$ appears \emph{inside the weight}, and the atoms \texttt{c(Y)} it sums over are guessed during search rather than given as input facts.
In general, an aggregate over guessed atoms may have to admit up to $2^{|Y|}$ potential values of the sum in the worst case (for instance, if the domain elements are powers of two, $Y \in \{1, 2, 4, 8, \ldots\}$).
Since the grounder cannot know what $S$ will be at decision time, it must enumerate every reachable sum. Even in this benchmark, where the domain $\{1,\ldots,n\}$ restricts the potential values to the dense interval $S \in \{0, \ldots, \frac{n(n+1)}{2}\}$, the grounder must still instantiate a separate ground directive for every pair $(X,S)$. This yields $\Theta(n^{3})$ ground objects (over one million, namely $1{,}010{,}200$, directives at $n = 100$) for the symmetric rules of $b/1$ and $c/1$. All of them instantiate the same single rule the modeller wrote, and all but the one matching the current partial sum are inert at any given decision.
A domain-specific heuristic introduced to help the engine solve a problem quicker ends up putting the answer out of reach.

\begin{figure}[H]
\centering
\includegraphics[width=\textwidth]{pipelines}
\caption{The same heuristic under the two approaches compared in this thesis.}
\label{fig:pipelines}
\end{figure}

The alternative is to leave the heuristic symbolic and to evaluate it while the search runs, when the partial assignment that $S$ refers to actually exists. This is the natural approach in lazy-grounding solvers such as Alpha \cite{weinzierl2017blending,comploi2023domain,comploi2023dynamicAggregates}. 

\section{Objective and Contributions}
\label{sec:contributions}

This work asks whether lazily evaluated heuristics can be added to clingo without giving up its ground-and-solve pipeline (Figure~\ref{fig:pipelines}), and whether the overhead of dynamic evaluation still allows the resulting system to outperform Alpha on the selected benchmarks.

The work contributes the following components and experimental analyses.

\begin{enumerate}
    \item \textbf{A lazy heuristic evaluator for clingo:} A propagator, registered through the public \texttt{clingo::Heuristic} interface, that keeps a declarative heuristic in memory as a compact description and evaluates it at every decision against the current partial assignment. 
    \item \textbf{Two propagator backends:} A compiled C++ backend, which materializes target candidates and maintains incremental aggregate states, and an in-process SWI-Prolog backend, which evaluates the heuristic as an ordinary logic program against the solver trail, read directly through foreign predicates.
    \item \textbf{Semantics and ablation study:} Ground-and-Solve, Lazy, clingo semantics and Alpha semantics combined give four variants of every encoding plus a heuristic-free baseline, to distinguish the effects of heuristic representation and interpretation.
    \item \textbf{Weight translation from Alpha's syntax:} Alpha's priority levels are global, while clingo's are local to a single atom. Section~\ref{sec:flattening} presents the transformation that folds a level into the weight and the conditions under which it preserves the intended order; Section~\ref{sec:weight-problem} the weight field limit imposed by clingo.
    \item \textbf{Trade-off analysis of the lazy approach:} By collecting per-decision metrics in both backends, we measure the runtime overhead of lazy heuristic evaluation against its savings in program size.
\end{enumerate}

The evaluation covers three problem families, the Balanced Sum Problem (BSP), the Partner Units Problem (PUP), and the House Reconfiguration Problem (HRP), on both backends. The implementation is publicly available.\footnote{\url{https://github.com/pierspad/clingo-lazy-heuristic-evaluator}}
